\clearpage
\phantomsection% 
\addcontentsline{toc}{chapter}{ABSTRACT}
\thispagestyle{fancy}

\begin{center}
	\large \bfseries \MakeUppercase{Abstract}\\
%	\normalsize \normalfont {\thetitleEN}\\
%	\normalsize \normalfont {\theauthor}\\
	\bigskip
	
	\normalsize \normalfont \justifying \singlespacing
	Class imbalance is a common problem in machine learning classification, causing models to be biased toward the majority class and fail to detect the minority class. HOUM (Hybrid Oversampling and Undersampling Method) is a hybrid method that combines Safe-Level SMOTE for oversampling and SVM for undersampling iteratively. However, previous research has not explored the effect of varying the regularization parameter $C$ and SVM kernel types, nor has it compared Safe-Level SMOTE with alternative oversampling techniques such as ADASYN. This study aims to analyze the effect of oversampling technique variation (Safe-Level SMOTE vs ADASYN) and SVM configuration (parameter $C$: 0.5; 1; 1.5; 2; 2.5; 3 and kernel types: Linear, RBF, Polynomial) within the HOUM method on imbalanced data classification performance. Four benchmark datasets were used, namely Diabetes, Haberman, Blood Transfusion, and Glass5, with varying imbalance ratios from 1:1.87 to 1:22.78. Data were split 80:20 using stratified split, balanced with 18 HOUM configuration combinations, and classified using Random Forest. Results show that HOUM-SLSMOTE consistently yields higher recall on Diabetes (0.6481), Haberman (0.3750), and Transfusion (0.7500), while HOUM-ADASYN excels in precision on Diabetes (0.7234) and Glass5 (0.6667). On Glass5 with extreme imbalance, HOUM-ADASYN achieves recall 1.0000, ROC-AUC 1.0000, and G-Mean 0.9877. The $C=0.5$ value is optimal for Diabetes, Haberman, and Glass5, whereas $C=2$ is more suitable for Transfusion. The RBF kernel performs best on the higher-dimensional dataset (Diabetes), the Linear kernel is more optimal on lower-dimensional datasets (Haberman, Transfusion), and on Glass5 all kernels yield identical metrics under HOUM-ADASYN. The best configurations are: HOUM-ADASYN with RBF $C=0.5$ (Diabetes, G-Mean 0.7401), HOUM-SLSMOTE with Linear $C=0.5$ (Haberman, G-Mean 0.4778), HOUM-SLSMOTE with Linear $C=2$ (Transfusion, G-Mean 0.6977), and HOUM-ADASYN with Linear $C=0.5$ (Glass5, G-Mean 0.9877). In conclusion, the selection of oversampling technique, SVM kernel, and $C$ value must be tailored to dataset characteristics due to the trade-off between recall and precision. HOUM is proven effective in improving minority class detection, including in extreme imbalance cases such as Glass5.
    
    \vspace{20pt}
	\raggedright\textbf{Keywords: HOUM, Oversampling, Undersampling, ADASYN, Safe-Level SMOTE, SVM, Imbalanced Data}
	
	\vfill
	
\end{center}
\clearpage